% Generated draft: requires mathematical and source review.
\documentclass[11pt]{article}
\usepackage[margin=25mm]{geometry}
\usepackage{fontspec}
\setmainfont{Latin Modern Roman}
\usepackage{amsmath,amssymb,mathtools,unicode-math}
\setmathfont{Latin Modern Math}
\usepackage{enumitem,needspace,xurl,hyperref}
\setlength{\parindent}{0pt}
\setlength{\parskip}{5pt}
\setlength{\emergencystretch}{3em}
\newcommand{\N}{\mathbb{N}}
\newcommand{\Z}{\mathbb{Z}}
\newcommand{\Q}{\mathbb{Q}}
\newcommand{\R}{\mathbb{R}}
\newcommand{\C}{\mathbb{C}}
\newcommand{\F}{\mathbb{F}}
\newcommand{\A}{\mathbb{A}}
\newcommand{\PP}{\mathbb{P}}
\newcommand{\E}{\mathbb{E}}
\newcommand{\eps}{\varepsilon}
\newcommand{\dd}{\,\mathrm{d}}
\begin{document}
\section*{UnsolvedMath ID: 30000087}
\textbf{Source problem code:} catalogue entry
\subsection{Definitions and mathematical statement}
The source requests an axiomatic graded structure, with a possibly ambiguous product, whose graded dimensions obey Macaulay's inequalities and, in exterior or squarefree settings, the Kruskal--Katona inequalities. It does not specify the axioms or a class of inputs on which the structure must exist.

For clarity, in the usual Hilbert-function setting write \(A=\bigoplus_{i\geq0}A_i\), with finite-dimensional pieces over a field \(k\), and \(h_i=\dim_k A_i\). For \(m>0\) and \(i\geq1\), write its unique binomial expansion
\[
m=\binom{a_i}{i}+\binom{a_{i-1}}{i-1}+\cdots+\binom{a_j}{j},
\qquad a_i>a_{i-1}>\cdots>a_j\geq j\geq1,
\]
and define
\[
m^{\langle i\rangle}=\binom{a_i+1}{i+1}+\cdots+\binom{a_j+1}{j+1},
\qquad 0^{\langle i\rangle}=0.
\]
The usual Macaulay growth inequalities are \(h_{i+1}\leq h_i^{\langle i\rangle}\). For standard graded commutative algebras, the usual setting includes \(A_0=k\) and generation by \(A_1\); the requested abstract structure is meant to justify comparable inequalities beyond an ordinary, unambiguously specified multiplication, not to drop all hypotheses on arbitrary graded dimensions.

The Kruskal--Katona target can be stated without an indexing ambiguity: for a family \(\mathcal F\) of \(i\)-element sets, its lower shadow is
\[
\partial\mathcal F=\{B:|B|=i-1,\ B\subset F\text{ for some }F\in\mathcal F\}.
\]
If \(|\mathcal F|=m\) has the expansion above, the shadow bound is
\[
|\partial\mathcal F|\geq\binom{a_i}{i-1}+\cdots+\binom{a_j}{j-1}.
\]
The source asks for an appropriate exterior or squarefree analogue within its proposed abstract framework, not for a new assertion about arbitrary sequences automatically satisfying these bounds.

The supplied assessment checked 2026-08-21 says that the extraction starts mid-discussion of an abstract graded ring and then turns to intersection homology; it contains no complete, uniquely testable problem. The meaning of ambiguous product, its compatibility with grading, the intended examples, and the role of intersection homology must therefore be recovered from the cited discussion before a precise existence or implication theorem can be formulated.
\subsection{Short English statement}
Seek axioms for a generalized graded structure whose dimensions satisfy Macaulay growth bounds, with Kruskal--Katona bounds in the exterior or squarefree setting. The numerical targets are recognizable, but the proposed structure and its axioms are missing.
\subsection{Sources}
\begin{itemize}
\item Local UnsolvedMath registry, permanent ID 30000087, catalogue entry, \emph{Abstract Graded Structures for Macaulay Inequalities}.
\item \emph{Combinatorial Commutative Algebra} (2005), Oberwolfach Reports, pp. 1703--1746, DOI \url{https://doi.org/10.4171/owr/2004/32}.
\item Kalai, supplied label \emph{Combinatorial Expectations}: \url{https://gilkalai.wordpress.com/wp-content/uploads/2019/03/ow04.pdf}.
\item Supplied literature triage dated 2026-08-21 identifies the extracted discussion as incomplete; no complete abstract structure or intersection-homology assertion is reproduced.
\end{itemize}
\end{document}
